% -*- latex -*-
%%%%%%%%%%%%%%%%%%%%%%%%%%%%%%%%%%%%%%%%%%%%%%%%%%%%%%%%%%%%%%%%
%%%%%%%%%%%%%%%%%%%%%%%%%%%%%%%%%%%%%%%%%%%%%%%%%%%%%%%%%%%%%%%%
%%%%
%%%% This text file is part of the lecture slides for
%%%% `Parallel Computing'
%%%% by Victor Eijkhout, copyright 2012-2026
%%%%
%%%% SPMD-slides.tex : slides about MPI's SPMD mode
%%%%
%%%%%%%%%%%%%%%%%%%%%%%%%%%%%%%%%%%%%%%%%%%%%%%%%%%%%%%%%%%%%%%%
%%%%%%%%%%%%%%%%%%%%%%%%%%%%%%%%%%%%%%%%%%%%%%%%%%%%%%%%%%%%%%%%

\begin{numberedframe}{Compiling}
  MPI compilers are usually called \indextermtt{mpicc},
  \indextermtt{mpif90}, \indextermtt{mpicxx}.

  These are not separate compilers,
  but scripts around the regular C/Fortran compiler. You can use all
  the usual flags.
\begin{lstlisting}[language=bash,numbers=none]
$ mpicc -show
icx -I/intel/include/stuff -L/intel/lib/stuff -Wwarnings # et cetera
\end{lstlisting}
(Python of course no compilation needed)

  (For CMake see slide~\ref{sl:mpi-cmake}.)
\end{numberedframe}

